\documentclass[12pt]{article}
\usepackage[margin=1in]{geometry}
\usepackage[T1]{fontenc}
\usepackage{lmodern}
\usepackage{microtype}
\usepackage{enumitem}
\usepackage[hidelinks]{hyperref}
\usepackage{xcolor}
\usepackage{titlesec}
\newcommand{\todo}[1]{\textcolor{red}{\textbf{#1}}}
\setlist{nosep,leftmargin=*}
\setlength{\parskip}{2pt}
\setlength{\parindent}{0pt}
\titlespacing*{\section}{0pt}{6pt}{2pt}
\titleformat*{\section}{\large\bfseries}

\title{\vspace{-2cm}Teaching BrickGPT New Bricks: Expanding a Brick Library In Context vs.\ In Weights\vspace{-0.5cm}}
\author{Team: \todo{[Name 1], [Name 2], [Name 3], [Name 4]} \quad (Topic 2: In-Context Learning)}
\date{\vspace{-1.2cm}}

\begin{document}
\maketitle

\section*{Abstract}
BrickGPT~\cite{pun2025brickgpt} generates stable, buildable brick structures from text, but only with eight small bricks (1$\times$1 to 2$\times$6).
We ask whether this specialist model can learn to use \emph{new} brick types (e.g.\ 4$\times$4, 2$\times$8) from its context alone, and how this compares with putting the new bricks into its weights.
We create new-brick data automatically by merging adjacent bricks in the existing StableText2Brick dataset, so no new labels are needed.
We compare zero-shot descriptions, few-shot demonstrations, LoRA, and a LoRA trained to ``read the brick library'' from context, tested on unseen brick types.
This is a concrete, checkable case of the course question: when should information live in the context, and when in the weights?

\section*{Introduction and Background}
\begin{itemize}
\item \textbf{BrickGPT}~\cite{pun2025brickgpt} fine-tunes Llama-3.2-1B-Instruct to predict the next brick as text on a 20$^3$ grid, with validity checks and physics-aware rollback; it names its ``fixed set of commonly used toy bricks'' as a limitation.
\item \textbf{Garg et al.}~\cite{garg2022can}: transformers learn simple function classes in context and match optimal estimators. \textbf{Brown et al.}~\cite{brown2020gpt3}: LLMs do new tasks from a few in-context examples.
\item \textbf{Wei et al.}~\cite{wei2023larger}: large models override semantic priors with in-context examples, but small and instruction-tuned models mostly cannot; BrickGPT is small and heavily fine-tuned.
\item \textbf{Agarwal et al.}~\cite{agarwal2024manyshot}: with hundreds of examples, ICL overrides pretraining biases and approaches fine-tuning.
\item \textbf{Lampinen et al.}~\cite{lampinen2025generalization}: with matched data, ICL generalizes more flexibly than fine-tuning; adding in-context inferences to fine-tuning data closes the gap.
\item \textbf{Chan et al.}~\cite{chan2022data}: ICL emerges from diverse, bursty data and competes with in-weights learning. \textbf{Singh et al.}~\cite{singh2023transient}: ICL can fade as in-weights solutions take over.
\item \textbf{Hu et al.}~\cite{hu2021lora}: LoRA, our in-weights method and BrickGPT's own training method. \textbf{Aky\"urek et al.}~\cite{akyurek2024ttt}: training on in-context examples at test time beats plain ICL on new tasks.
\end{itemize}

\section*{Key Questions}
\begin{enumerate}
\item \textbf{Zero/few-shot:} Given a description of new bricks (and optionally $k$ demonstrations), does BrickGPT use them, and do structures stay valid and stable?
\item \textbf{Context vs.\ weights:} At matched data, how does in-context use compare with LoRA on the same examples, in usage, quality and cost?
\item \textbf{Generalization:} If trained to read a brick library from its context, can the model use brick types never seen in training? Plain LoRA cannot, by construction.
\end{enumerate}

\section*{Hypothesis}
Zero-shot description barely changes behavior: the fine-tuned prior for the 8 bricks dominates~\cite{wei2023larger}. Few-shot raises usage but lowers validity. Plain LoRA is best on the bricks it was trained on but does not transfer. The library-in-context LoRA approaches plain LoRA on seen bricks and clearly beats all others on held-out bricks~\cite{lampinen2025generalization}.

\section*{Methods}
\textbf{Task design.} Any shape can be built from 1$\times$1 bricks, so large bricks must help: the prompt gives a \emph{brick budget} (``at most $N$ bricks'', $N$ below the 8-brick solution), and brick count and stability are metrics, since large bricks bridge seams with fewer pieces.

\textbf{Data (no new labels).} From StableText2Brick (47k structures, 28k ShapeNet objects, GPT-4o captions), greedily merge adjacent aligned bricks into larger types (e.g.\ 4$\times$4, 2$\times$8, 4$\times$6). Shape and caption are unchanged; we re-check stability and keep stable structures. New types are split into \emph{seen} (train) and \emph{held-out} (test only). We extend the brick library in BrickGPT's validity checker, stability analysis and rendering.

\textbf{Conditions} (BrickGPT checkpoint): (A) original 8-brick model; (B) zero-shot description; (C) few-shot, $k\in\{1,2,5,10\}$; (D) LoRA on $n$ merged examples, $n$ matched to C and scaled to 100--5k; (E) library-in-context LoRA with randomly sampled libraries, tested on held-out bricks.

\textbf{Metrics.} BrickGPT's \% valid, \% stable, mean/min stability, CLIP, DINO; plus new-brick usage rate, brick count, budget satisfaction, voxel IoU vs.\ ground truth. Train/val loss curves for D and E; error bars over 3 seeds and 100+ test prompts.

\textbf{Compute.} BrickGPT reports 3 epochs on 47k structures in about 12 h on 8 RTX A6000s. Our LoRA runs use 100--5k examples on a 1B model, so a few GPU-hours each, about \todo{[X]} GPU-hours in total; inference with rollback is the main cost. Available: \todo{[GPU type, count, hours/week --- TO BE FILLED]}.

\textbf{Code bases.} BrickGPT repo (MIT; TRL/Accelerate fine-tuning, stability analysis, rendering), Hugging Face PEFT.

\begin{thebibliography}{99}\setlength{\itemsep}{-3pt}\footnotesize
\bibitem{pun2025brickgpt} A.~Pun, K.~Deng, R.~Liu, D.~Ramanan, C.~Liu, J.-Y.~Zhu. Generating Physically Stable and Buildable Brick Structures from Text. ICCV 2025 (Marr Prize).
\bibitem{garg2022can} S.~Garg et al. What Can Transformers Learn In-Context? A Case Study of Simple Function Classes. NeurIPS 2022.
\bibitem{brown2020gpt3} T.~Brown et al. Language Models are Few-Shot Learners. NeurIPS 2020.
\bibitem{wei2023larger} J.~Wei et al. Larger Language Models Do In-Context Learning Differently. arXiv:2303.03846, 2023.
\bibitem{agarwal2024manyshot} R.~Agarwal et al. Many-Shot In-Context Learning. NeurIPS 2024.
\bibitem{lampinen2025generalization} A.~K.~Lampinen et al. On the Generalization of Language Models from In-Context Learning and Finetuning. arXiv:2505.00661, 2025.
\bibitem{chan2022data} S.~C.~Y.~Chan et al. Data Distributional Properties Drive Emergent In-Context Learning in Transformers. NeurIPS 2022.
\bibitem{singh2023transient} A.~K.~Singh et al. The Transient Nature of Emergent In-Context Learning in Transformers. NeurIPS 2023.
\bibitem{hu2021lora} E.~J.~Hu et al. LoRA: Low-Rank Adaptation of Large Language Models. ICLR 2022.
\bibitem{akyurek2024ttt} E.~Aky\"urek et al. The Surprising Effectiveness of Test-Time Training for Few-Shot Learning. arXiv, 2024.
\end{thebibliography}

\end{document}
